% ~600 words

\subsection{Data Collection Pipeline}

The dataset was constructed through a nine-stage pipeline
(\Cref{tab:data_pipeline} and \Cref{fig:data_pipeline_pub}), implemented in the open-source
\emph{ConnectedLecturer} project.\footnote{\url{https://github.com/TUBAF-IFI-ConnectedLecturer/Data_aggregation}}
Starting from GitHub's Search API, we identified \num{1076} repositories
containing potential LiaScript content.
A tree-level aggregation step extracted \num{57096} Markdown files,
of which \num{3317} (\SI{5.8}{\percent}) were validated as LiaScript courses
using a combination of rule-based filters and an LLM classifier
(Llama~3.3 70B).
For each course, the pipeline collected commit history (\num{19366} commits),
parsed header metadata (\SI{99.4}{\percent} success rate),
extracted \num{44} binary feature flags, and performed AI-based enrichment
(titles, keywords, educational level) for \SI{94.6}{\percent} of courses.
The final consolidated dataset comprises \num{3317} courses with
166~columns each (collected 2026-03-13).

%\begin{table}[ht]
%  \centering
%  \caption{Data collection pipeline stages}
%  \label{tab:pipeline}
%  \begin{tabular}{llr}
%    \toprule
%    \textbf{Phase} & \textbf{Method} & \textbf{$n$} \\
%    \midrule
%    Repository search    & GitHub Search API           & \num{1076} repos \\
%    File extraction      & Tree API + heuristics       & \num{57096} files \\
%    Validation           & Rules + LLM                 & \num{3317} courses (\SI{5.8}{\percent}) \\
%    Commit analysis      & Commits API                 & \num{19366} commits \\
%    Header metadata      & Regex parsing               & \num{3297} courses (\SI{99.4}{\percent}) \\
%    Feature detection    & 44 regex patterns            & $\num{3317} \times 44$ \\
%    AI enrichment        & RAG + Llama 3.3 70B         & \num{3138} courses (\SI{94.6}{\percent}) \\
%    User profiles        & GitHub Users API             & 445 profiles \\
%    Consolidation        & Inner join                   & $\num{3317} \times 166$ \\
%    \bottomrule
%  \end{tabular}
%\end{table}

\begin{table}[t]
\centering
\caption{Nine-stage data construction pipeline for the ConnectedLecturer dataset.}
\label{tab:data_pipeline}
\begin{tabular}{p{2.2cm} p{3.2cm} p{2.6cm} p{2.8cm}}
\toprule
\textbf{Phase} & \textbf{Method} & \textbf{Output} & \textbf{Coverage} \\
\midrule
Repository search 
& GitHub Search API 
& 1{,}076 repositories 
& -- \\

File extraction 
& Tree API + heuristics 
& 57{,}096 Markdown files 
& -- \\

Validation 
& Rule-based filters + LLM classifier 
& 3{,}317 LiaScript courses 
& 5.8\% of files \\

Commit analysis 
& Commits API 
& 19{,}366 commits 
& -- \\

Header metadata 
& Regex parsing 
& Metadata for 3{,}297 courses 
& 99.4\% of courses \\

Feature detection 
& 44 regex patterns 
& 44 features for 3{,}317 courses 
& 100\% \\

AI enrichment 
& RAG + Llama 3.3 70B 
& Enrichment for 3{,}138 courses 
& 94.6\% of courses \\

User profiles 
& GitHub Users API 
& 445 profiles 
& -- \\

Consolidation 
& Course-centered join 
& Final dataset: 3{,}317 courses $\times$ 166 columns 
& -- \\
\bottomrule
\end{tabular}
\end{table}

\begin{figure}[htbp]
\centering
\begin{tikzpicture}[
    font=\small,
    node distance=8mm and 12mm,
    >={Stealth[length=3mm,width=2mm]},
    base/.style={
        rectangle,
        rounded corners=2.2mm,
        draw,
        line width=0.8pt,
        align=center,
        minimum height=1.35cm,
        text width=4.0cm,
        inner sep=4pt
    },
    search/.style={
        base,
        draw=blue!55!black,
        fill=blue!8
    },
    validate/.style={
        base,
        draw=orange!70!black,
        fill=orange!12
    },
    enrich/.style={
        base,
        draw=teal!60!black,
        fill=teal!10
    },
    aux/.style={
        base,
        draw=purple!55!black,
        fill=purple!9
    },
    finalbox/.style={
        base,
        draw=black,
        fill=gray!12,
        text width=5.1cm,
        minimum height=1.5cm
    },
    flow/.style={
        ->,
        line width=1.0pt,
        draw=black!75
    },
    groupbox/.style={
        rounded corners=2.5mm,
        draw=black!45,
        dashed,
        line width=0.7pt,
        inner sep=7pt
    },
    grouplabel/.style={
        font=\footnotesize\bfseries,
        fill=white,
        inner sep=2pt
    }
]

% --- Main linear path --------------------------------------------------------

\node[search] (search) {Repository search\\
GitHub Search API\\[1mm]
\textbf{1,076 repositories}};

\node[search, below=of search] (files) {File extraction\\
Tree API + heuristics\\[1mm]
\textbf{57,096 Markdown files}};

\node[validate, below=15mm of files] (valid) {Validation\\
Rules + LLM\\[1mm]
\textbf{3,317 LiaScript courses}\\
\emph{(5.8\% of files)}};

% --- Parallel branches -------------------------------------------------------

\node[enrich, below left=37mm and -18mm of valid] (commits) {Commit analysis\\
Commits API\\[1mm]
\textbf{19,366 commits}};

\node[enrich, below=14mm of valid] (header) {Header metadata\\
Regex parsing\\[1mm]
metadata for \textbf{3,297 courses}\\
\emph{(99.4\% coverage)}};

\node[enrich, below right=37mm and -18mm of valid] (features) {Feature detection\\
44 regex patterns\\[1mm]
\textbf{44 features} for\\
\textbf{3,317 courses}};

\node[enrich, above right= 6mm and -18mm of features] (ai) {AI enrichment\\
RAG + Llama 3.3 70B\\[1mm]
enrichment for \textbf{3,138 courses}\\
\emph{(94.6\% coverage)}};

\node[aux, above left= 8mm and -18mm of commits] (users) {User profiles\\
GitHub Users API\\[1mm]
\textbf{445 profiles}};

% --- Final node --------------------------------------------------------------

\node[finalbox, below=40mm of header] (final) {Course-centered consolidation\\
validated courses combined with derived metadata\\[1mm]
\textbf{Final dataset: 3,317 courses $\times$ 166 columns}};

% --- Arrows ------------------------------------------------------------------

\draw[flow] (search) -- (files);
\draw[flow] (files) -- (valid);

\draw[flow] (valid.190) -| (commits);
\draw[flow] (valid) -- (header);
\draw[flow] (valid.350) -| (features);
\draw[flow] (valid.east) -| (ai.north);
\draw[flow] (valid.west) -| (users.north);

\draw[flow] (users.south) |- (final.west);
\draw[flow] (commits.south) -- (final.150);
\draw[flow] (header) -- (final);
\draw[flow] (features.south) -- (final.30);
\draw[flow] (ai.south) |- (final.east);

% --- Group annotations -------------------------------------------------------

\begin{scope}[on background layer]
    \node[groupbox, fit=(search)(files)] (g1) {};
    \node[grouplabel, above=1mm of g1.north] {Candidate identification};

    \node[groupbox, fit=(valid)] (g2) {};
    \node[grouplabel, above=1mm of g2.north] {Course validation};

    \node[groupbox, fit=(users)(commits)(header)(features)(ai)] (g3) {};
    \node[grouplabel, above=1mm of g3.north] {Parallel analysis and enrichment};

    \node[groupbox, fit=(final)] (g4) {};
    \node[grouplabel, above=1mm of g4.north] {Final assembly};
\end{scope}

\end{tikzpicture}
\caption{Publication-oriented overview of the ConnectedLecturer data construction workflow. After candidate identification and course validation, several analysis and enrichment branches are executed in parallel and merged into the final consolidated course dataset.}
\label{fig:data_pipeline_pub}
\end{figure}

Our primary unit of analysis is the individual \emph{course file}
(one validated \texttt{.md} file = one course).
For grouping purposes courses are segmented by their \emph{repository
context} (Section~\ref{sec:three_groups}).

\subsection{Feature Taxonomy}
\label{sec:taxonomy}

For each course file the pipeline extracts \num{44} binary feature flags
via syntactic pattern matching on the Markdown source.\footnote{All regex patterns and their mapping to feature labels are documented at \url{https://github.com/TUBAF-IFI-ConnectedLecturer/Data_aggregation/blob/main/cl-pipeline/run/pipelines/liascript/docs/LiaScript_Features_Patterns.md}}
We organise these features into four didactically motivated categories
(Table~\ref{tab:taxonomy}).

\begin{table}[ht]
  \centering
  \caption{LiaScript feature taxonomy (4 categories, 44 tracked features)}
  \label{tab:taxonomy}
  \begin{tabular}{p{2.8cm}p{3.2cm}p{5.2cm}}
    \toprule
    \textbf{Category} & \textbf{Design intent} & \textbf{Representative features} \\
    \midrule
    \textcolor{catPresentation}{\textbf{Presentation}} &
      Self-paced, accessible delivery &
      Narrator (TTS), animations, effects, galleries \\[4pt]
    \textcolor{catInteraction}{\textbf{Interaction}} &
      Active learning \& formative feedback &
      Single/multiple choice, text input, matrix quiz, surveys \\[4pt]
    \textcolor{catReuse}{\textbf{Reuse}} &
      Scalable authoring via shared components &
      Imports, custom macros, official templates \\[4pt]
    \textcolor{catEmbedding}{\textbf{Embedding}} &
      Hands-on executable environments &
      Executable code, WebApps, external scripts, ASCII diagrams \\
    \bottomrule
  \end{tabular}
\end{table}

Each category addresses a different pedagogical barrier.
\emph{Presentation} features target accessibility (screen-reader-friendly TTS,
animated step-by-step reveals).
\emph{Interaction} features enable formative self-assessment without a learning
management system.
\emph{Reuse} features support the Don't-Repeat-Yourself principle at course level,
allowing shared building blocks across a curriculum.
\emph{Embedding} features close the gap between explanation and practice by
running code, simulations, or visualisations directly in the browser.

\subsection{Three-Group Segmentation}
\label{sec:three_groups}

Each course resides in a GitHub repository owned by an individual or
organisation account.
The \num{3317} courses were contributed by \num{307} unique committers
working across \num{205} repository accounts (\num{167}~individual,
\num{37}~organisational, 1~unclassified).
Organisation accounts act as hosting containers; the actual authors are
identified through the Git commit history (\texttt{contributors\_list}).

To prevent distortions from authors with fundamentally different
relationships to LiaScript, we segment courses into three groups based
on the \emph{repository context}:
\begin{description}
  \item[Internal] (462~courses, 7~repository accounts, 58~committers):
    Repositories owned by the core development team and their
    organisational accounts
    (SebastianZug, andre-dietrich, LiaPlayground, LiaScript, LiaBooks,
    LiaTemplates, TUBAF-IfI-LiaScript).
    These include templates, demonstrations, and the developers' own
    teaching materials.
  \item[MINT-the-GAP] (1\,147~courses, 1~organisation account, 3~committers):
    A single prolific school-focused initiative producing structured
    STEM courses.
    With nearly a third of all courses, this account would dominate any
    community-level analysis if not separated.
  \item[Community] (1\,708~courses, 197~repository accounts, 269~committers):
    All remaining courses, representing independent adoption by
    educators without direct involvement in LiaScript development.
    The group spans a wide range of engagement:
    82~accounts have produced exactly one course --- including
    test repositories, workshop materials, and one-off teaching
    projects --- while 8~accounts each contribute 56--146~courses.
    Despite this spread, no single account dominates
    (the largest accounts for \SI{8.5}{\percent}),
    and feature adoption rates show no systematic divergence between
    prolific and occasional contributors.
    Excluding all single-course accounts shifts adoption rates by
    less than one percentage point, confirming that the group is not
    distorted by test or tutorial content.
\end{description}

\noindent
Because grouping follows the repository context, the same committer
may appear in multiple groups.
In particular, 21~committers --- including the two core developers ---
contribute to both Internal and Community repositories.
Their work in Community repositories is counted as Community,
reflecting the collaborative context rather than the contributor's
identity.
For \num{306} courses (\SI{9.2}{\percent}) no commit history was
available; these are assigned solely based on their repository account.

\subsection{Complementary User Survey}
\label{sec:survey}

To triangulate the corpus findings, we conducted a structured needs
assessment with seven active LiaScript authors from different institutions
(March~2026).
Participants prioritised feature requests and usability issues on a
five-point scale.
The resulting 60~items are analysed qualitatively in
Section~\ref{sec:implications}.

\subsection{Limitations}

The corpus is restricted to \emph{public} GitHub repositories.
Courses hosted on other platforms, behind authentication, or in private
repositories are not captured.
Feature detection via pattern matching may miss features used through
indirect mechanisms (e.g.\ macros that generate quizzes programmatically).
The three-group segmentation assigns courses by repository context,
so the 21~committers active in multiple groups are counted in each
group separately rather than as a single entity.
The AI enrichment step (titles, keywords, DDC classification) was not
used in the present analysis but is available for future work.
